% =====================================================================
%  1. PROBLEM STATEMENT  -  Người viết: TV5
%  Đề yêu cầu: vấn đề đặt ra, tầm quan trọng, ứng dụng.
%  Tự viết bằng lời của mình. Mọi ý lấy từ nguồn ngoài phải \cite{...}.
%  Độ dài gợi ý: 1 - 1,5 trang.
% =====================================================================
\section{Problem statement}\label{sec:problem}\writer{TV5}

\subsection{The problem}
% Gợi ý: ảnh số là gì (lưới pixel, mỗi pixel 3 số R, G, B trong 0..255);
% bài toán: biến đổi màu cả ảnh bằng cách nhân ma trận 3x3 với vector màu
% của từng pixel; câu hỏi chính: phép lọc nào đảo ngược được, vì sao.
\todo{mô tả bài toán}

\subsection{Why it matters}
% Gợi ý: một ví dụ trực quan về phép nhân ma trận, định thức, nghịch đảo,
% số điều kiện; liên hệ giữa "khả nghịch về lý thuyết" và "khôi phục được thật".
\todo{tầm quan trọng}

\subsection{Applications}
% Gợi ý: bộ lọc ảnh trên điện thoại / mạng xã hội, chuyển không gian màu,
% chỉnh cân bằng trắng, hỗ trợ người mù màu, nén ảnh... (có trích dẫn).
\todo{ứng dụng}

\subsection{Objectives and structure of the report}
\todo{mục tiêu nhóm đặt ra và tóm tắt các phần sau}
